\section{Parameter Sweep}
The preferred open-ended analysis is a gap sweep over representative bus-ring gaps of \SI{80}{\nano\metre}, \SI{100}{\nano\metre}, \SI{120}{\nano\metre}, and \SI{150}{\nano\metre}. No numerical sweep data was found in \path{results/}, so the section is presented as a qualitative design discussion and a TODO table.

\begin{table}[H]
\centering
\caption{Gap sweep template. Numerical entries require exported FDE, INTERCONNECT, or varFDTD sweep data.}
\label{tab:gap-sweep}
\begin{tabular}{lllll}
\toprule
Gap (nm) & \(|k|^2\) & FSR (nm) & \(Q\) & Drop-port comment \\
\midrule
80 & TODO & TODO & TODO & Not available from current simulation output \\
100 & TODO & TODO & TODO & Not available from current simulation output \\
120 & TODO & TODO & TODO & Not available from current simulation output \\
150 & TODO & TODO & TODO & Not available from current simulation output \\
\bottomrule
\end{tabular}
\end{table}

The expected physical trend is that increasing the gap reduces evanescent coupling, so \(|k|\) decreases and \(|t|\) increases. Weaker coupling generally narrows the resonance linewidth and can increase \(Q\) and finesse, but it may also reduce drop-port efficiency and make the sensor response more sensitive to fabrication error and loss. Conversely, a smaller gap increases coupling and drop-port transfer but can broaden the resonance.
